\documentclass[11pt]{amsart}

\usepackage[T1]{fontenc}
\usepackage{lmodern}
\usepackage{mathtools,amssymb,amsthm}
\usepackage{microtype}
\usepackage[hidelinks]{hyperref}

\newtheorem{theorem}{Theorem}[section]
\newtheorem{proposition}[theorem]{Proposition}
\newtheorem{lemma}[theorem]{Lemma}
\newtheorem{corollary}[theorem]{Corollary}
\theoremstyle{definition}
\newtheorem{definition}[theorem]{Definition}
\newtheorem{algorithm}[theorem]{Algorithm}
\theoremstyle{remark}
\newtheorem{remark}[theorem]{Remark}
\newtheorem{example}[theorem]{Example}

\newcommand{\Z}{\mathbb Z}
\newcommand{\Bn}{\mathcal B_n}
\newcommand{\An}{\mathcal A_n}
\newcommand{\Bodd}{\mathcal B_n^{\mathrm{odd}}}
\newcommand{\Beven}{\mathcal B_n^{\mathrm{even}}}
\newcommand{\Boddres}{\mathcal B_{n,\circ}^{\mathrm{odd}}}
\newcommand{\Ares}{\mathcal A_{n,\circ}}
\newcommand{\Sn}{\mathcal S_n}
\newcommand{\Tn}{\mathcal T_n}
\newcommand{\Sres}{\mathcal S_n^{\circ}}
\newcommand{\Tres}{\mathcal T_n^{\circ}}
\newcommand{\Sbar}{\overline{\mathcal S}_n}
\newcommand{\Tbar}{\overline{\mathcal T}_n}
\newcommand{\Dcal}{\mathcal P_Q}
\newcommand{\key}{\operatorname{key}}
\newcommand{\pc}{\operatorname{pc}}
\newcommand{\sgn}{\operatorname{sgn}}
\newcommand{\cont}{\operatorname{cont}}
\newcommand{\Span}{\operatorname{span}_{\Z}}
\newcommand{\Qip}[2]{\left\langle #1,#2\right\rangle_Q}
\newcommand{\ceil}[1]{\left\lceil #1\right\rceil}
\newcommand{\floor}[1]{\left\lfloor #1\right\rfloor}

\title[A two-to-one map for Tunnell's forms]
{Stable Matching and a Two-to-One Map
for Tunnell's Ternary Forms}

\date{August 2026}

\subjclass[2020]{Primary 05C70; Secondary 11E25, 11D85, 68W40}
\keywords{globally ranked pairs, stable matching, parking function, ternary quadratic form, arithmetic correspondence}
\hypersetup{
  pdftitle={Stable Matching and a Two-to-One Map for Tunnell's Ternary Forms},
  pdfauthor={},
  pdfsubject={An arithmetic map constructed by stable matching},
  pdfkeywords={globally ranked pairs, stable matching, parking function, ternary quadratic form, arithmetic correspondence}
}

\begin{document}

\begin{abstract}
For every odd positive squarefree congruent number $n$, we construct a
deterministic map $\Phi_n:\Bn\to\An$ whose fibers have cardinality two, where
$\Bn$ and $\An$ are the representation sets of $n$ by
$2x^2+y^2+8z^2$ and $2x^2+y^2+32z^2$.  Tunnell's theorem supplies the
identity $|\Bn|=2|\An|$; more generally, the construction applies
to every odd positive squarefree $n$ satisfying this identity.  Halving the
third coordinate gives the even branch.  On the odd branch, three integral
quarter-turns define a partial bijection.  We complete it by passing the
residual points to their antipodal orbits and applying stable matching, with
preferences determined by primitive midpoint directions.  Equal-rank edges
are disjoint, so the stable matching is unique.  Its preference lists can be
generated arithmetically from rank-two lattice cosets, and the proposal counts
form a parking function with total at most $m(m+1)/2$, where $m$ is the number
of residual orbits on either side.
\end{abstract}

\maketitle

\section{Introduction}

Tunnell in \cite{Tunnell} related the congruent-number problem to
representation numbers of ternary quadratic forms.  In this paper, we use
his counting identity to construct a two-to-one map between the corresponding
representation sets.  For an odd positive integer $n$, we set
\begin{align}
 \Bn&=\{(x,y,z)\in\Z^3:2x^2+y^2+8z^2=n\},\label{eq:B-def}\\
 \An&=\{(u,v,w)\in\Z^3:2u^2+v^2+32w^2=n\}.\label{eq:A-def}
\end{align}
A ternary quadratic form is a homogeneous quadratic polynomial in three
variables, and a representation of $n$ is simply an integer triple at which
the form has value $n$.  We call $n$ \emph{squarefree} if no prime square
divides it.  Both sets above are finite because each squared coordinate is
bounded by $n$.

A positive integer is \emph{congruent} if it is the area of a right triangle
with rational side lengths.  Equivalently, the elliptic curve
$y^2=x^3-n^2x$ has a rational point of infinite order; see
\cite[Chap.~I, \S9]{Koblitz} and \cite[pp.~323--324]{Tunnell}.  A consequence
of Tunnell's theorem is that every odd positive squarefree congruent number
satisfies
\begin{equation}\label{eq:Tunnell-balance}
 |\Bn|=2|\An|.
\end{equation}
We use this identity from Tunnell's theorem.  The construction below does not
require any further modular-form or elliptic-curve arguments.

\begin{example}[The balance for $n=41$]\label{ex:balance-41}
For $n=41$, a solution in $\mathcal B_{41}$ must satisfy
$|x|\le4$, $|y|\le6$, and $|z|\le2$; a solution in $\mathcal A_{41}$ must
satisfy the same first two bounds and $|w|\le1$.  Sorting by the last
coordinate gives
\[
\begin{array}{c|ccccc|c}
\mathcal B_{41}:z&0&1&-1&2&-2&\text{total}\\ \hline
\text{number of }(x,y)&4&8&8&6&6&32
\end{array}
\]
and
\[
\begin{array}{c|ccc|c}
\mathcal A_{41}:w&0&1&-1&\text{total}\\ \hline
\text{number of }(u,v)&4&6&6&16.
\end{array}
\]
Thus $|\mathcal B_{41}|=2|\mathcal A_{41}|$.  For example,
\[
 (2,-1,2)\in\mathcal B_{41},
 \qquad
 (2,-1,1)\in\mathcal A_{41},
\]
and the even-branch map simply halves the last coordinate.
\end{example}

Bach and Ryan \cite{BachRyan} give methods for computing the representation
numbers in this identity.  Here we take \eqref{eq:Tunnell-balance} as given
and ask how to pair the individual representations.  We will work with any
odd positive squarefree $n$ satisfying this identity, so the construction is
not restricted to the congruent-number case.  We do not reprove Tunnell's
identity or give a new congruent-number criterion.

Separate $\Bn$ according to the parity of its third coordinate:
\begin{align*}
 \Beven&=\{(x,y,z)\in\Bn:z\equiv0\pmod2\},\\
 \Bodd&=\{(x,y,z)\in\Bn:z\equiv1\pmod2\}.
\end{align*}
The map
\begin{equation}\label{eq:even-branch}
 (x,y,2w)\longmapsto(x,y,w)
\end{equation}
is a bijection $\Beven\xrightarrow{\sim}\An$.  Consequently
\eqref{eq:Tunnell-balance} is equivalent to $|\Bodd|=|\An|$, and the
essential problem is to construct a bijection from the odd branch to the
target set.

Three quarter-turns give a bijection on part of $\Bodd$.  We call the points
outside its domain and image the residual points.  To match them, we pass
to their $\{\pm1\}$-orbits and join every source orbit to every target orbit.
A primitive midpoint direction assigns a rank to each edge.  Both endpoints
order their incident edges by this rank, while equal-rank edges are
vertex-disjoint.  This is an instance of stable marriage with globally ranked
pairs in the sense of Abraham, Levavi, Manlove, and O'Malley
\cite{AbrahamEtAl}.  We will show that the stable matching is unique.

Stable matching and parking functions are classical objects
\cite{GaleShapley,GusfieldIrving,KonheimWeiss,StanleyParking}.  In our setting,
the midpoint directions determine the preferences, and the numbers of
proposals form a parking function.  The counting identity alone would allow
an arbitrary pairing.  Our map is instead determined by the quarter-turns,
the midpoint directions, and the orientation conventions fixed below.  It
commutes with $X\mapsto-X$, and the same matching determines its inverse.

\begin{theorem}[A two-to-one map]\label{thm:main}
Let $n$ be an odd positive squarefree integer satisfying
$|\Bn|=2|\An|$.  Fix the orientation and ordering conventions introduced
below.  Then the construction defines a deterministic map
\[
 \Phi_n:\Bn\longrightarrow\An
\]
such that every element of $\An$ has exactly two preimages.  On $\Beven$, the
map is \eqref{eq:even-branch}.  On $\Bodd$, it is the disjoint union of three
quarter-turn branches and the unique stable matching of the remaining orbit
sets.  The algorithms below compute this matching from arithmetically
generated preference lists, without first storing its complete edge table.
\end{theorem}

Here a \emph{fiber} is the set of all points mapping to one target.  The map
comes from two bijections.  The first gives the even preimage of each target,
and the second gives the odd preimage:
\begin{equation}\label{eq:master-map}
\begin{gathered}
 \Beven\xrightarrow[\text{bijection}]{\ (x,y,2w)\mapsto(x,y,w)\ }\An,
 \\[4pt]
 \Bodd=D_n\sqcup\Boddres
 \xrightarrow[\text{bijection}]{\ \Lambda_n\,\sqcup\,\Psi_n^{\circ}\ }
 I_n\sqcup\Ares=\An.
\end{gathered}
\end{equation}
The residual bijection is obtained by changing coordinates, matching
antipodal orbits, and lifting back to points:
\begin{equation}\label{eq:residual-flow}
 \Boddres\xrightarrow{\ \iota_B\ }\Sres
 \longrightarrow\Sbar
 \xrightarrow{\ \overline M_n\ }\Tbar
 \longrightarrow\Tres
 \xrightarrow{\ \iota_A^{-1}\ }\Ares.
\end{equation}
The unlabeled arrows pass to antipodal orbits and then lift back with a
specified sign.  The cardinality hypothesis is used to prove that
$\overline M_n$ is perfect; it is not used to choose the edge ranks.

\begin{theorem}[Congruent-number case]\label{thm:congruent}
For every odd positive squarefree congruent number $n$, the construction
defines a deterministic map $\Phi_n:\Bn\to\An$ such that every element of
$\An$ has exactly two preimages.
\end{theorem}

Section~\ref{sec:coordinates} introduces the two coordinate models, and
Section~\ref{sec:quarter-turn} constructs the direct branches.
Section~\ref{sec:projective} constructs the residual matching, and
Section~\ref{sec:assembly} defines the two-to-one map.
Sections~\ref{sec:affine-stream}--\ref{sec:example} explain how to compute the
matching and work out the residual example at $n=41$.
Section~\ref{sec:complexity} bounds the number of proposals and the cost of
deferred acceptance.  We finish with two questions about the residual map.

We will use two examples at $n=41$.  The first uses a quarter-turn, and the
second uses the residual matching:
\[
\begin{array}{c|c|c|c}
\text{case}&\text{odd source}&\text{target}&\text{even source}\\ \hline
\text{direct quarter-turn}&(-4,-1,1)&(2,-1,1)&(2,-1,2)\\
\text{residual matching}&(-2,-5,1)&(4,-3,0)&(4,-3,0).
\end{array}
\]
Each row lists the two preimages of one target.  The quarter-turns and the
residual matching apply to disjoint subsets of $\Bodd$.

\section{Two changes of coordinates}
\label{sec:coordinates}

We now introduce two changes of coordinates.  The first is used for the
quarter-turns and the second for midpoint directions.  Each is applied to
the original representation sets:
\[
\begin{array}{ccccc}
\Bodd&\xleftrightarrow[\ e_B^{-1}\ ]{\ e_B\ }&L_1(n)
&\xrightarrow{\ J_1,J_2,J_3\ }&L_0(n)
\xleftrightarrow[\ e_A\ ]{\ e_A^{-1}\ }\An,\\[5pt]
\Bodd&\xleftrightarrow[\ \iota_B^{-1}\ ]{\ \iota_B\ }&\Sn
&\xleftrightarrow[\text{midpoint lines}]{}&\Tn
\xleftrightarrow[\ \iota_A\ ]{\ \iota_A^{-1}\ }\An.
\end{array}
\]
In the first row, both sets are represented by sums of three squares, where
we can apply coordinate rotations.  In the second, both lie on the same
quadratic surface, and their half-sums and half-differences are integral.

The following table records the notation used below:
\begin{center}
\renewcommand{\arraystretch}{1.12}
\begin{tabular}{@{}c p{0.70\textwidth}@{}}
$D_n,I_n,\Lambda_n$&direct odd-source domain, its target image, and the
quarter-turn bijection between them;\\
$\Boddres,\Ares,\Psi_n^\circ$&points left after the quarter-turn bijection and their
residual bijection;\\
$\Sn,\Tn$&midpoint-coordinate lifts of the odd source and target sets;\\
$\Sbar,\Tbar$&antipodal residual orbit sets;\\
$\Dcal$&canonically oriented primitive midpoint directions;\\
$\Psi_n,\Phi_n$&the full odd-branch bijection and the final two-to-one map.
\end{tabular}
\end{center}

\subsection{Sum-of-three-squares coordinates for the direct map}
For $\epsilon\in\{0,1\}$, put
\[
 \begin{aligned}
 L_\epsilon(n)=\{(a,b,c)\in\Z^3:
 &\ a^2+b^2+c^2=n,\\
 &\ a\equiv b\pmod4,\quad
 \frac{a-b}{4}\equiv\epsilon\pmod2\}.
 \end{aligned}
\]

\begin{lemma}[Sum-of-three-squares coordinates]\label{lem:pure-coordinates}
The maps
\[
 e_B(x,y,z)=(x+2z,x-2z,y),
 \qquad
 e_A(u,v,w)=(u+4w,u-4w,v)
\]
are bijections
\[
 e_B:\Bodd\xrightarrow{\sim}L_1(n),
 \qquad
 e_A:\An\xrightarrow{\sim}L_0(n),
\]
with inverses
\[
 e_B^{-1}(a,b,c)=\left(\frac{a+b}{2},c,\frac{a-b}{4}\right),
\]
\[
 e_A^{-1}(a,b,c)=\left(\frac{a+b}{2},c,\frac{a-b}{8}\right).
\]
\end{lemma}

\begin{proof}
The norm identities are immediate:
\[
 (x+2z)^2+(x-2z)^2+y^2=2x^2+y^2+8z^2,
\]
\[
 (u+4w)^2+(u-4w)^2+v^2=2u^2+v^2+32w^2.
\]
The congruence conditions record whether $(a-b)/4$ is odd or even.
Conversely, if $(a,b,c)\in L_1(n)$, then $(a+b)/2$ and $(a-b)/4$ are
integral, the latter is odd, and the first norm identity gives a point of
$\Bodd$.  If $(a,b,c)\in L_0(n)$, then $(a-b)/8$ is integral and the second
identity gives a point of $\An$.  The displayed inverse formulas therefore
give inverses on both sides.
\end{proof}

\begin{example}[The two coordinate changes at $n=41$]
\label{ex:coordinates-41}
Consider the odd source $X=(-4,-1,1)\in\mathcal B_{41}^{\mathrm{odd}}$.
Lemma~\ref{lem:pure-coordinates} sends it to
\[
 e_B(X)=(-2,-6,-1),
 \qquad
 (-2)^2+(-6)^2+(-1)^2=41.
\]
Here $(-2)-(-6)=4$, so $(a-b)/4=1$ is odd, as required for $L_1(41)$.
Likewise, for $Y=(2,-1,1)\in\mathcal A_{41}$,
\[
 e_A(Y)=(6,-2,-1),
 \qquad
 6^2+(-2)^2+(-1)^2=41,
\]
and $(6-(-2))/4=2$ is even, as required for $L_0(41)$.  The inverse formulas
recover $X$ and $Y$ from these triples.
\end{example}

\subsection{Midpoint coordinates for the residual graph}
Put
\[
 Q(X_1,X_2,X_3)=2X_1^2+X_2^2+2X_3^2
\]
and write
\[
 \Qip{X}{Y}=2X_1Y_1+X_2Y_2+2X_3Y_3
\]
for its associated bilinear form.  Thus
$Q(X+Y)=Q(X)+2\Qip{X}{Y}+Q(Y)$.  Define
\begin{align*}
 \Sn&=\{E\in\Z^3:Q(E)=n,\ E_3\equiv2\pmod4\},\\
 \Tn&=\{F\in\Z^3:Q(F)=n,\ F_3\equiv0\pmod4\}.
\end{align*}
The lifts
\[
 \iota_B(x,y,z)=(x,y,2z),
 \qquad
 \iota_A(u,v,w)=(u,v,4w)
\]
identify $\Bodd$ with $\Sn$ and $\An$ with $\Tn$.

\begin{lemma}[Midpoint parity]\label{lem:midpoint-parity}
If $E\in\Sn$ and $F\in\Tn$, then
\[
 \frac{E+F}{2},\qquad \frac{E-F}{2}
\]
are integral vectors with odd third coordinate.
\end{lemma}

\begin{proof}
Since $n$ is odd, the second coordinates of both $E$ and $F$ are odd.  Their
third coordinates are congruent to $2$ and $0$ modulo $4$, respectively.
Consequently, modulo $8$ one has
\[
 n\equiv2E_1^2+1\equiv2F_1^2+1,
\]
so $E_1$ and $F_1$ have the same parity.  The second-coordinate halves are
integral because both entries are odd, and the third-coordinate halves are
odd because $E_3\equiv2\pmod4$ and $F_3\equiv0\pmod4$.  Thus both displayed
vectors are integral and have odd third coordinate.
\end{proof}

For $0\ne v\in\Z^3$, define
\[
 \cont(v)=\gcd(|v_1|,|v_2|,|v_3|).
\]
Call $v$ \emph{primitive} if $\cont(v)=1$.  Let $\Dcal$ contain the unique
representative, with first nonzero coordinate positive, of every pair
$\{d,-d\}$ of primitive vectors with $d_3$ odd.  Order $\Dcal$
lexicographically by
\begin{equation}\label{eq:direction-key}
 \key(d)=\bigl(Q(d),d_1,d_2,d_3\bigr).
\end{equation}
For directions we choose the sign with first nonzero coordinate positive.
For antipodal point orbits, introduced below, we instead choose the
lexicographically smaller point.

\begin{lemma}[Primitive generator of the midpoint line]
\label{lem:midpoint-direction}
For every $E\in\Sn$ and $F\in\Tn$, there are unique $d\in\Dcal$ and a
nonzero odd integer $q$ such that
\[
 \frac{E+F}{2}=qd.
\]
They satisfy
\[
 \Qip{E}{d}=qQ(d),
 \qquad
 F=-E+2qd.
\]
For a fixed $d$, the edges $E\leftrightarrow F$ obtained in this way are
pairwise disjoint.
\end{lemma}

\begin{proof}
By Lemma~\ref{lem:midpoint-parity}, the midpoint is a nonzero integral vector
with odd third coordinate.  Dividing by its content and applying the
orientation convention gives $q$ and $d$, with $q$ odd.  If
$h=(E-F)/2$, equality of the norms gives $\Qip{d}{h}=0$, and hence
$\Qip{E}{d}=qQ(d)$ and $F=-E+2qd$.  Uniqueness follows from the unique
sign-normalized primitive generator of
$\mathbb Q(E+F)\cap\mathbb Z^3$, the integer points on the rational line
spanned by $E+F$.  For fixed $E$ and $d$, the identity
$q=\Qip{E}{d}/Q(d)$ determines $q$, and then the formula determines $F$.
Interchanging the roles of the endpoints gives the same conclusion at a
target.  Hence the edges belonging to a fixed direction are pairwise
disjoint.
\end{proof}

\begin{example}[A midpoint direction at $n=41$]
\label{ex:midpoint-41}
Take the lifted residual points
\[
 E=(-2,-5,2)\in\mathcal S_{41},
 \qquad
 F=(4,-3,0)\in\mathcal T_{41}.
\]
Both have $Q$-value $41$.  Their half-sum and half-difference are
\[
 \frac{E+F}{2}=(1,-4,1),
 \qquad
 \frac{E-F}{2}=(-3,-1,1).
\]
Both vectors are integral and have odd third coordinate, as in
Lemma~\ref{lem:midpoint-parity}.  The first vector is already
primitive and canonically oriented, hence $q=1$ and
$d=(1,-4,1)$.  Moreover,
\[
 Q(d)=20,
 \qquad
 \Qip{E}{d}=20=qQ(d),
 \qquad
 -E+2d=F.
\]
The other midpoint line is generated, after sign normalization, by
$e=(3,1,-1)$, whose key begins with $Q(e)=21$.  Thus the key of $d$ is the
smaller of the two.  When we pass to antipodal orbits, this comparison will
determine which sign of the target is used.
\end{example}

\section{A direct quarter-turn correspondence}\label{sec:quarter-turn}

Define three integral rotations in the sum-of-three-squares model by
\[
 J_1(a,b,c)=(a,-c,b),
 \qquad
 J_2(a,b,c)=(c,b,-a),
 \qquad
 J_3(a,b,c)=(-b,a,c).
\]
Each is an element of $\operatorname{SO}_3(\Z)$: it is an integer matrix that
preserves squared length and orientation.  Geometrically, $J_i$ fixes one
coordinate axis and rotates the other two coordinates through a right angle.

\begin{proposition}[Quarter-turn domains]\label{prop:quarter-turn-domains}
Let $(a,b,c)\in L_1(n)$.  Then
\begin{align*}
 J_1(a,b,c)\in L_0(n)&\iff a+c\equiv0\pmod8,\\
 J_2(a,b,c)\in L_0(n)&\iff c-b\equiv0\pmod8,\\
 J_3(a,b,c)\in L_0(n)&\iff a+b\equiv0\pmod8.
\end{align*}
At most one of these three congruences holds.
\end{proposition}

\begin{proof}
The differences between the first two coordinates of the three images are
$a+c$, $c-b$, and $-(a+b)$, respectively.  Membership in $L_0(n)$ is
therefore equivalent to the stated divisibility by $8$.

In original odd-source coordinates, $c=y$ is odd and
$a-b=4z\equiv4\pmod8$.  If any two of the displayed congruences held, then,
after eliminating $a$ or $b$, one would obtain
\[
 a-b\equiv -2c\pmod8.
\]
Since $c$ is odd, the right-hand side is $2$ or $6$ modulo $8$, contradicting
$a-b\equiv4\pmod8$.  Hence the three domains are pairwise disjoint.
\end{proof}

Writing the rotations in the original coordinates gives the
following partial map.  For $(x,y,z)\in\Bodd$ with
$x+2z+y\equiv0\pmod8$, set
\begin{equation}\label{eq:Lambda1}
 \Lambda_1(x,y,z)=
 \left(\frac{x+2z-y}{2},\ x-2z,\ \frac{x+2z+y}{8}\right).
\end{equation}
If $y-x+2z\equiv0\pmod8$, set
\begin{equation}\label{eq:Lambda2}
 \Lambda_2(x,y,z)=
 \left(\frac{x-2z+y}{2},\ -x-2z,\ \frac{y-x+2z}{8}\right).
\end{equation}
Finally, if $x\equiv0\pmod4$, set
\begin{equation}\label{eq:Lambda3}
 \Lambda_3(x,y,z)=\left(2z,\ y,\ -\frac{x}{4}\right).
\end{equation}
The domains are pairwise disjoint by Proposition~\ref{prop:quarter-turn-domains}.

\begin{theorem}[Direct partial bijection]\label{thm:direct-layer}
The formulas \eqref{eq:Lambda1}--\eqref{eq:Lambda3} define a partial
bijection
\[
 \Lambda_n:D_n\xrightarrow{\sim}I_n
\]
from a subset $D_n\subseteq\Bodd$ to a subset $I_n\subseteq\An$.  Its inverse
is determined from a target $(u,v,w)\in\An$ by the following mutually
exclusive conditions and formulas:
\begin{align*}
 u+4w-v\equiv4\pmod8:
 &\quad
 \left(\frac{u+4w+v}{2},-u+4w,\frac{u+4w-v}{4}\right),\\
 u-4w+v\equiv4\pmod8:
 &\quad
 \left(\frac{u-4w-v}{2},u+4w,\frac{4w-u-v}{4}\right),\\
 u\equiv2\pmod4:
 &\quad
 \left(-4w,v,\frac{u}{2}\right).
\end{align*}
Moreover, $\Lambda_n$ is equivariant under the antipodal involution
$X\mapsto-X$.
\end{theorem}

\begin{proof}
The construction is $e_A^{-1}J_i e_B$ on each branch, so norm preservation and
integrality follow from Proposition~\ref{prop:quarter-turn-domains}.  The
inverse formulas are $e_B^{-1}J_i^{-1}e_A$.  Their conditions state exactly
that the inverse image lies in $L_1(n)$.  They are mutually exclusive: the first
two cannot both hold because $v$ is odd, and either of them is incompatible
with $u\equiv2\pmod4$ by parity.  Direct composition gives the two identity
maps.  Negation equivariance follows because each $J_i$ is linear and each
domain condition is unchanged by negation.
\end{proof}

\begin{example}[The three quarter-turn branches]
\label{ex:quarter-turn-41}
The first two branches already occur at $n=11$; the direct construction does
not require the cardinality hypothesis.  Direct substitution gives
\[
\begin{array}{c|c|c|c}
\text{branch}&X&e_B(X)&\Lambda_i(X)\\ \hline
\Lambda_1&(-1,-1,1)&(1,-3,-1)&(1,-3,0)\\
\Lambda_2&(-1,1,-1)&(-3,1,1)&(1,3,0).
\end{array}
\]
For the first row, $a+c=1-1=0$, so $J_1$ applies and
$J_1(1,-3,-1)=(1,1,-3)=e_A(1,-3,0)$.  For the second row,
$c-b=1-1=0$, so $J_2$ applies and
$J_2(-3,1,1)=(1,1,3)=e_A(1,3,0)$.

For $n=41$, return to $X=(-4,-1,1)$ from
Example~\ref{ex:coordinates-41}.  Its
sum-of-three-squares coordinates are $(a,b,c)=(-2,-6,-1)$.  Since
\[
 a+b=-8\equiv0\pmod8,
\]
Proposition~\ref{prop:quarter-turn-domains} selects $J_3$ and neither of the
other two quarter-turns.  Indeed,
\[
 J_3(-2,-6,-1)=(6,-2,-1)=e_A(2,-1,1).
\]
In the original coordinates this is exactly
\[
 \Lambda_3(-4,-1,1)=(2,-1,1).
\]
The inverse condition is $u\equiv2\pmod4$, and the third inverse formula gives
$(-4w,v,u/2)=(-4,-1,1)$.
\end{example}

Set
\[
 \Boddres=\Bodd\setminus D_n,
 \qquad
 \Ares=\An\setminus I_n,
\]
and let $\Sres=\iota_B(\Boddres)$ and
$\Tres=\iota_A(\Ares)$.
Under the hypothesis of Theorem~\ref{thm:main},
\begin{equation}\label{eq:residual-balance}
 |\Sres|=|\Tres|.
\end{equation}
Indeed, $|\Bodd|=|\An|$ and $\Lambda_n$ removes equally many points from the
two sides.

\section{Stable matching on antipodal orbits}
\label{sec:projective}

The residual sets are invariant under the antipodal involution $X\mapsto-X$.
Since $n>0$, no point in either set equals its negative, so each orbit has
two elements.  Define the orbit sets
\[
 \Sbar=\Sres/\{\pm1\},
 \qquad
 \Tbar=\Tres/\{\pm1\}.
\]
By \eqref{eq:residual-balance}, these sets have the same cardinality; write
\[
 |\Sbar|=|\Tbar|=m.
\]
Each orbit is an antipodal pair $[X]=\{X,-X\}$.  Choose the lexicographically
smaller of $X$ and $-X$ as its canonical representative $\pc([X])$.  Let
\[
 G_n=K_{\Sbar,\Tbar}
\]
be the complete bipartite graph on the two orbit sets.

Fix $[E]\in\Sbar$ and $[F]\in\Tbar$, with canonical representatives $E$ and
$F$.  Apply Lemma~\ref{lem:midpoint-direction} to the two edges between
chosen representatives
$E\leftrightarrow F$ and $E\leftrightarrow -F$.  Let their direction keys be
$k_+(E,F)$ and $k_-(E,F)$.

\begin{lemma}[Edge ranking on antipodal orbits]\label{lem:projective-key}
The two keys $k_+(E,F)$ and $k_-(E,F)$ are distinct.  The quantity
\[
 \kappa([E],[F])=\min\{k_+(E,F),k_-(E,F)\}
\]
is independent of the chosen representatives.  There is a unique sign
\[
 \eta([E],[F])\in\{\pm1\}
\]
such that the edge between representatives
\[
 E\longleftrightarrow \eta([E],[F])F
\]
has key $\kappa([E],[F])$.
\end{lemma}

\begin{proof}
Write
\[
 \frac{E+F}{2}=qd,
 \qquad
 \frac{E-F}{2}=ce
\]
with primitive oriented directions $d,e$.  Equality of the norms gives
$\Qip{d}{e}=0$.  The vectors are nonzero, so positive definiteness implies
$d\ne e$ and hence the keys are distinct.  Replacing either representative by
its negative interchanges the two signed midpoint lines, so their unordered
pair of keys is unchanged.  The smaller one determines the unique sign.
\end{proof}

\begin{example}[Choosing the sign of a target]
\label{ex:signed-orbit-41}
Take the canonical representatives
\[
 E=(-2,-5,-2),
 \qquad
 F=(-4,-3,0)
\]
of the orbits later denoted by $s_1$ and $t_1$.  We compare the edges from
$E$ to $F$ and to $-F$.  For $E\leftrightarrow F$, the midpoint is
\[
 \frac{E+F}{2}=(-3,-4,-1)=-(3,4,1),
\]
so its oriented primitive direction has key $(36,3,4,1)$.  For
$E\leftrightarrow -F$, the midpoint is
\[
 \frac{E-F}{2}=(1,-1,-1),
\]
whose key is $(5,1,-1,-1)$.  The second key is smaller.  Therefore
\[
 \kappa([E],[F])=(5,1,-1,-1),
 \qquad
 \eta([E],[F])=-1,
\]
and the corresponding map on representatives is
\[
 E\longmapsto-F=(4,3,0).
\]
For the other source point $-E$, multiplying both endpoints by $-1$ gives the
other lift $-E\mapsto F$.  The two source points are therefore paired with
the two target points, and these pairings can be read in either direction.
\end{example}

If the actual source is $\sigma E$, where $\sigma\in\{\pm1\}$, the lift of an
orbit edge is
\begin{equation}\label{eq:sign-lift}
 \sigma E\longmapsto \sigma\eta([E],[F])F.
\end{equation}
The same rule, read from the target, recovers the source.

\begin{proposition}[Equal-key classes are matchings]
\label{prop:key-class-matching}
For a fixed key $k$, the edges of $G_n$ with key $k$ form a matching: no two
share a source or target orbit.
\end{proposition}

\begin{proof}
At a fixed signed endpoint, a primitive direction determines at most one
opposite signed endpoint by Lemma~\ref{lem:midpoint-direction}.  It remains to
check that the two signs of an endpoint do not create two orbits.  If a fixed
direction sends $F$ to $E=-F+2qd$, then at the endpoint $-F$ the multiplier is
$-q$ and the partner is
\[
 -(-F)+2(-q)d=-E.
\]
Thus changing the sign of one endpoint changes the sign of the other and does
not change its orbit.  Hence a fixed key occurs at most once at either
endpoint of $G_n$.
\end{proof}

\begin{example}[Two disjoint edges with the same key]
\label{ex:equal-key-41}
The key
\[
 (5,1,-1,-1)
\]
occurs on the edge $(s_1,t_1)$ computed in
Example~\ref{ex:signed-orbit-41}.  For a second edge, take
\[
 E'=(-2,5,-2)=\pc(s_3),
 \qquad
 F'=(0,-3,4)=\pc(t_4).
\]
Their midpoint is $(-1,1,1)=-(1,-1,-1)$, so this edge has the same key.
The sources $s_1,s_3$ and targets $t_1,t_4$ are distinct.  Hence the two
edges with the same key are disjoint.
\end{example}

Every vertex ranks its incident edges by increasing $\kappa$.
Proposition~\ref{prop:key-class-matching} implies that equal-key edges are
vertex-disjoint, so every vertex has a strict preference order.  A pair
$(s,t)\notin M$ \emph{blocks} a matching $M$ if both $s$ and $t$ prefer each
other to their partners in $M$, with being unmatched least preferred.  A
matching with no blocking pair is \emph{stable}.  Thus $G_n$ is a
stable-marriage instance with globally ranked pairs
\cite{GaleShapley,GusfieldIrving,AbrahamEtAl}: the two preference systems are
induced by one global ranking of edges.

\begin{proposition}[Unique stable matching]\label{prop:unique-stable}
The graph $G_n$ has a unique stable matching.  It is obtained by processing
the edge keys in increasing order and selecting every edge whose two
endpoints are still unmatched.  Since $|\Sbar|=|\Tbar|$, this stable matching
is perfect.
\end{proposition}

\begin{proof}
First run the stated greedy procedure.  If an edge is not selected, then at
least one of its endpoints was already matched when that edge was processed.
The selected edge at that endpoint has smaller key; equality is impossible
by Proposition~\ref{prop:key-class-matching}.  Thus an unselected edge cannot
block, and the greedy matching is stable.  It is also perfect: otherwise
balance would leave an unmatched source and an unmatched target, but the edge
joining them would have been selected when it was processed.

It remains to prove uniqueness.  Consider the least key among the edges whose
endpoints have not already been removed.  Every edge of that key is forced in
any stable matching: neither endpoint has a smaller edge in the remaining
graph, so omission would make the two endpoints a blocking pair.
Proposition~\ref{prop:key-class-matching} says that these edges have disjoint
endpoints, and hence they can be selected and removed simultaneously.
A removed endpoint prefers its selected edge to every edge processed later,
and the restriction of a stable matching to the remaining vertices is stable.
Induction therefore shows that every stable matching agrees with the greedy
matching.
\end{proof}

\begin{example}[A two-by-two globally ranked instance]
\label{ex:ranked-2x2}
Before returning to the arithmetic graph, consider two sources
$s_1,s_2$ and two targets $t_1,t_2$.  Give the four edges the global ranks
\[
\begin{array}{c|cc}
 &t_1&t_2\\ \hline
s_1&1&3\\
s_2&2&4
\end{array}
\]
with smaller ranks preferred.  The greedy rule first selects $(s_1,t_1)$.
The edges of ranks $2$ and $3$ then meet an already matched vertex, so the
next selected edge is $(s_2,t_2)$.  Hence
\[
 M=\{(s_1,t_1),(s_2,t_2)\}.
\]
The only other perfect matching is not stable, because $(s_1,t_1)$ blocks it.
This is the simplest instance of Proposition~\ref{prop:unique-stable}.
\end{example}

\section{The two-to-one map}\label{sec:assembly}

Let
\[
 \overline M_n:\Sbar\xrightarrow{\sim}\Tbar
\]
be the unique stable matching of Proposition~\ref{prop:unique-stable}.  Lift it
to residual points using \eqref{eq:sign-lift}, and return to the original
coordinates.  Denote the resulting residual bijection by
\[
 \Psi_n^{\circ}:\Boddres\xrightarrow{\sim}\Ares.
\]
Define the odd-branch bijection
\[
 \Psi_n(X)=
 \begin{cases}
  \Lambda_n(X),&X\in D_n,\\
  \Psi_n^{\circ}(X),&X\in\Boddres.
 \end{cases}
\]

\begin{proof}[Proof of Theorem~\ref{thm:main}]
The map $\Lambda_n$ is a bijection $D_n\to I_n$ by
Theorem~\ref{thm:direct-layer}.  The residual orbit matching is perfect by
Proposition~\ref{prop:unique-stable}, and the sign lift pairs the two points in
each source orbit with the two points in the matched target orbit.  Hence
$\Psi_n$ is a bijection $\Bodd\to\An$.

Define
\[
 \Phi_n(x,y,z)=
 \begin{cases}
  (x,y,z/2),&z\text{ even},\\
  \Psi_n(x,y,z),&z\text{ odd}.
 \end{cases}
\]
For $(u,v,w)\in\An$, the even branch contributes the preimage $(u,v,2w)$,
and the odd branch contributes the unique preimage $\Psi_n^{-1}(u,v,w)$.
They are distinct because their third coordinates have opposite parity, and
there are no further preimages on either branch.  Thus every fiber has
cardinality two.
\end{proof}

\begin{proof}[Proof of Theorem~\ref{thm:congruent}]
For an odd squarefree congruent number $n$, Tunnell's theorem gives
\eqref{eq:Tunnell-balance}, so Theorem~\ref{thm:main} applies.
\end{proof}

\begin{remark}[Role of cardinality]\label{rem:cardinality}
The formulas defining the quarter-turns and edge keys make sense without
\eqref{eq:Tunnell-balance}.  The equality is used only to make the two residual
orbit sets equally large; Proposition~\ref{prop:unique-stable} then makes their
stable matching perfect.  The later generator and deferred-acceptance
procedure compute this already defined matching.
\end{remark}

\begin{example}[The direct $n=41$ fiber]
\label{ex:fiber-41}
Let $Y=(2,-1,1)\in\mathcal A_{41}$.  The even branch gives
\[
 (2,-1,2)\longmapsto(2,-1,1),
\]
while Example~\ref{ex:quarter-turn-41} gives the odd preimage
\[
 (-4,-1,1)\longmapsto(2,-1,1).
\]
Their last coordinates have opposite parity, so
\[
 \Phi_{41}^{-1}(2,-1,1)
 =\{(2,-1,2),(-4,-1,1)\}.
\]
These are the two preimages listed in the first row of the introductory table.
\end{example}

\section{Arithmetic generation of preference lists}\label{sec:affine-stream}

We now explain how to compute the
residual matching without first constructing all $m^2$ edges.  For a fixed
source orbit, its preference list can be recovered by searching for the
primitive midpoint directions of its incident edges.  We first bound this
search and then write it as a problem in a rank-two affine lattice.

\begin{lemma}[Bound for the preferred direction]
\label{lem:projective-half-bound}
If $d$ is the direction realizing the smaller signed orientation of a
residual orbit edge, then
\[
 Q(d)\le\frac{n-1}{2}.
\]
\end{lemma}

\begin{proof}
The two midpoint decompositions from Lemma~\ref{lem:projective-key} give
\[
 E=qd+ce,
 \qquad
 F=qd-ce.
\]
Since $d$ and $e$ are orthogonal, we also have
\[
 n=q^2Q(d)+c^2Q(e).
\]
The chosen orientation gives $\key(d)<\key(e)$, hence $Q(d)\le Q(e)$.  Since
$q$ and $c$ are nonzero,
\[
 n\ge Q(d)+Q(e)\ge2Q(d).
\]
The claim follows because $n$ is odd.
\end{proof}

\begin{example}[The midpoint bound is attained]
In Example~\ref{ex:midpoint-41}, the preferred direction satisfies
\[
 Q(d)=20=\frac{41-1}{2}.
\]
Thus the bound in Lemma~\ref{lem:projective-half-bound} can be an equality.
\end{example}

\subsection{An orthogonal-complement lattice}

We use the isometric embedding
\[
 \tau(X_1,X_2,X_3)=(X_1+X_3,X_2,X_1-X_3).
\]
Then
\[
 \|\tau(X)\|^2=Q(X),
 \qquad
 \tau(X)\cdot\tau(Y)=\Qip{X}{Y}.
\]
The image of $\tau$ is the index-two lattice
\begin{equation}\label{eq:tau-lattice}
 \mathcal L_\tau=\{r\in\Z^3:r_1\equiv r_3\pmod2\},
 \qquad
 \tau^{-1}(r)=
 \left(\frac{r_1+r_3}{2},r_2,\frac{r_1-r_3}{2}\right).
\end{equation}
Let us fix a residual source orbit.  We write $E$ for its canonical lifted
representative and set $p=\tau(E)$.  Since $\|p\|^2=n$ and $n$ is squarefree,
$p$ is primitive; otherwise the square of a common divisor of its coordinates
would divide $n$.

We first describe which lattice points are needed.
If $h=Q(d)$ and $t=qh$, then $r=\tau(d)$ is an integer point in
\begin{equation}\label{eq:plane-sphere}
 \underbrace{\{r:p\cdot r=t\}}_{\text{an affine plane}}
 \ \cap\ 
 \underbrace{\{r:\|r\|^2=h\}}_{\text{a sphere}}
 \ \cap\ \mathcal L_\tau.
\end{equation}
The integer points parallel to this plane form the rank-two lattice
$p^\perp\cap\Z^3$.  We next choose coordinates on this lattice and use them to
turn the plane-sphere intersection into a quadratic equation in two integers.

\begin{lemma}[A basis of the orthogonal lattice]
\label{lem:orthogonal-frame}
There are deterministically computable vectors $e_1,e_2,z_0\in\Z^3$ such that
\[
 e_1\times e_2=p,
 \qquad
 p\cdot z_0=1.
\]
The pair $e_1,e_2$ is an oriented basis of $p^\perp\cap\Z^3$ and may be
Lagrange--Gauss reduced while preserving $e_1\times e_2=p$.
\end{lemma}

\begin{proof}
Since $p$ is primitive, there is a B\'ezout vector $z_0$ with
$p\cdot z_0=1$.  Applying Smith normal form to the row vector $p$ gives an
integral basis $e_1,e_2$ of its kernel.  Every $x\in\Z^3$ has the unique
decomposition
\[
 x=\bigl(x-(p\cdot x)z_0\bigr)+(p\cdot x)z_0.
\]
Thus $e_1,e_2,z_0$ form a basis of $\Z^3$.  The vector $e_1\times e_2$ is an
integral multiple of the primitive normal vector $p$.  Taking its scalar
product with $z_0$ shows that the multiple is $\pm1$, and reversing the kernel
basis if necessary gives $e_1\times e_2=p$.  The usual integral
Lagrange--Gauss steps
\[
 (e_1,e_2)\mapsto(e_1,e_2-me_1),
 \qquad
 (e_1,e_2)\mapsto(e_2,-e_1)
\]
have determinant $1$, preserve the cross product, and terminate with
$|2e_1\cdot e_2|\le\|e_1\|^2\le\|e_2\|^2$.
\end{proof}

We use Smith normal form in its usual algorithmic form: integer column
operations, obtained from the Euclidean algorithm, change the primitive row
$p$ to $(1,0,0)$.  Transforming the last two standard basis vectors back gives
a basis of $p^\perp\cap\Z^3$.  We take nonnegative Euclidean remainders and use
lexicographic tie breaking.  For Lagrange--Gauss reduction, we round ties in
the nearest-integer step toward zero and choose the lexicographically smaller
oriented basis when the two reduced vectors have the same norm.  These choices
make the basis and B\'ezout vector deterministic.  They do not affect the final
preference list, which is sorted by the intrinsic direction key.  We refer to
\cite[Chap.~2]{Cohen} for these standard lattice algorithms.

We write
\[
 \alpha=\|e_1\|^2,
 \quad
 \beta=e_1\cdot e_2,
 \quad
 \gamma=\|e_2\|^2,
\]
\[
 \delta=z_0\cdot e_1,
 \quad
 \varepsilon=z_0\cdot e_2,
 \quad
 \zeta=\|z_0\|^2.
\]
The identity $e_1\times e_2=p$ gives
\begin{equation}\label{eq:gram-det}
 \alpha\gamma-\beta^2=n.
\end{equation}
This follows from
\[
 \alpha\gamma-\beta^2
 =\|e_1\times e_2\|^2
 =\|p\|^2=n.
\]

\subsection{Affine norm equations}

For an edge incident to $E$, write its midpoint data as
\[
 \frac{E+F}{2}=qd,
 \qquad
 h=Q(d),
 \qquad
 r=\tau(d).
\]
We then have
\begin{equation}\label{eq:Pr}
 \|r\|^2=h,
 \qquad
 p\cdot r=qh.
\end{equation}
We set $t=qh$.  The frame from the previous subsection now gives coordinates
on the affine plane $p\cdot r=t$.

\begin{lemma}[Parametrization of a lattice coset]\label{lem:affine-coset}
Every integral solution of $p\cdot r=t$ is uniquely of the form
\[
 r=tz_0+\mu e_1+\nu e_2,
 \qquad \mu,\nu\in\Z.
\]
For fixed $h,q,\nu$, the condition $\|r\|^2=h$ is equivalent to
\begin{equation}\label{eq:quadratic-mu}
 \alpha\mu^2+2(\beta\nu+t\delta)\mu
 +\gamma\nu^2+2t\varepsilon\nu+t^2\zeta-h=0.
\end{equation}
\end{lemma}

\begin{proof}
The difference $r-tz_0$ is orthogonal to $p$, and $e_1,e_2$ form an integral
basis of the orthogonal-complement lattice.  Expanding the squared norm gives
\eqref{eq:quadratic-mu}.
\end{proof}

To determine the possible values of $\nu$, we write
\[
 K=\beta\delta-\alpha\varepsilon,
 \qquad
 M=h(n-q^2h).
\]
The discriminant of \eqref{eq:quadratic-mu}, after dividing by $4$, is
\[
 \Delta_\nu=(\beta\nu+t\delta)^2-
 \alpha(\gamma\nu^2+2t\varepsilon\nu+t^2\zeta-h).
\]

\begin{lemma}[Bounds for the affine coordinates]\label{lem:short-interval}
One has
\begin{equation}\label{eq:short-identity}
 n\Delta_\nu=\alpha M-(n\nu-tK)^2.
\end{equation}
Put $R=\floor{\sqrt{\alpha M}}$.  By \eqref{eq:short-identity},
$\Delta_\nu\ge0$ exactly when $\nu$ lies in the interval
\begin{equation}\label{eq:nu-interval}
 \ceil{\frac{tK-R}{n}}
 \le \nu\le
 \floor{\frac{tK+R}{n}}.
\end{equation}
For each such $\nu$, let $s=\floor{\sqrt{\Delta_\nu}}$.  If
$s^2=\Delta_\nu$, the possible values of $\mu$ form the set
\begin{equation}\label{eq:mu-roots}
 \left\{
 \frac{-(\beta\nu+t\delta)+u}{\alpha}:
 u\in\{s,-s\},\quad
 \alpha\mid-(\beta\nu+t\delta)+u
 \right\}.
\end{equation}
If $s^2\ne\Delta_\nu$, there are no roots.  In particular, when
$\Delta_\nu=0$, the set \eqref{eq:mu-roots} contains at most one integer.
\end{lemma}

\begin{proof}
The Gram matrix of the basis $e_1,e_2,z_0$ is
\[
 G=\begin{pmatrix}
 \alpha&\beta&\delta\\
 \beta&\gamma&\varepsilon\\
 \delta&\varepsilon&\zeta
 \end{pmatrix}.
\]
Since $\det[e_1\ e_2\ z_0]=1$, one has $\det G=1$.  Expanding gives
\[
 n\zeta-\bigl(\alpha\varepsilon^2-2\beta\delta\varepsilon
 +\gamma\delta^2\bigr)=1,
\]
and a direct rearrangement using $n=\alpha\gamma-\beta^2$ gives
\[
 K^2+n(\delta^2-\alpha\zeta)=-\alpha.
\]
The discriminant can now be written as
\[
 \Delta_\nu=-n\nu^2+2tK\nu+t^2(\delta^2-\alpha\zeta)+\alpha h.
\]
After multiplying by $n$ and using the preceding identity together with
$t=qh$, we obtain
\[
 n\Delta_\nu
 =\alpha(nh-t^2)-(n\nu-tK)^2
 =\alpha h(n-q^2h)-(n\nu-tK)^2,
\]
which is \eqref{eq:short-identity}.  Since $n\nu-tK$ is integral, its
square is at most $\alpha M$ precisely when
$|n\nu-tK|\le\floor{\sqrt{\alpha M}}$; this is equivalent to
\eqref{eq:nu-interval}.  The quadratic formula gives
\eqref{eq:mu-roots}, interpreted as a set.
\end{proof}

\begin{example}[Recovering one $n=41$ edge arithmetically]
\label{ex:lattice-41}
Return to $E=(-2,-5,2)$ and $F=(4,-3,0)$ from
Example~\ref{ex:midpoint-41}.  The isometry gives
\[
 p=\tau(E)=(0,-5,-4).
\]
One convenient oriented kernel basis and B\'ezout vector are
\[
 e_1=(1,0,0),
 \qquad
 e_2=(0,-4,5),
 \qquad
 z_0=(0,-1,1).
\]
These vectors follow directly from the two defining equations.  The kernel
equation $p\cdot r=0$ is $-5r_2-4r_3=0$, so $r_2=-4k$ and $r_3=5k$, while
$r_1$ is free.  This gives $e_1,e_2$, and the B\'ezout identity $5-4=1$
gives $z_0$.  Hence
$e_1\times e_2=p$ and $p\cdot z_0=1$.  For this frame,
\[
 \begin{aligned}
  \alpha&=1,& \beta&=0,& \gamma&=41,\\
  \delta&=0,& \varepsilon&=9,& \zeta&=2,& K&=-9.
 \end{aligned}
\]

The midpoint direction is $d=(1,-4,1)$, so
\[
 h=Q(d)=20,\qquad q=1,\qquad t=20,\qquad M=20(41-20)=420.
\]
Since $R=\floor{\sqrt{420}}=20$ and $tK=-180$, the interval
\eqref{eq:nu-interval} becomes
\[
 \ceil{\frac{-200}{41}}\le \nu\le
 \floor{\frac{-160}{41}},
\]
which leaves only $\nu=-4$.  The discriminant identity then gives
\[
 41\Delta_{-4}=420-\bigl(41(-4)-20(-9)\bigr)^2=164,
 \qquad \Delta_{-4}=4.
\]
Thus \eqref{eq:mu-roots} gives $\mu=\pm2$.  The choice $\mu=2$ produces
\[
 r=20z_0+2e_1-4e_2=(2,-4,0),
 \qquad
 \tau^{-1}(r)=(1,-4,1)=d.
\]
The other root has negative first coordinate and is removed by the orientation
convention.  The outer loop also uses negative $q$, so the canonically oriented
opposite direction is recovered there whenever it gives the preferred signed
edge.  Finally,
\[
 -E+2d=(4,-3,0)=F.
\]
Since the canonical representative of $[F]$ is $F_0=(-4,3,0)$, the generator
records
\[
 \bigl((20,1,-4,1),[F_0],-1\bigr).
\]
This recovers one entry of the preference list directly from the arithmetic
of the source $E$.
\end{example}

\subsection{Generating a preference list}

The preceding calculation gives a finite procedure for producing the complete
preference list of a fixed source orbit.

\begin{algorithm}[Generation of a preference list]\label{alg:incident-stream}
Fix a residual source orbit $[E]$ and let $E=\pc([E])$.  We perform the
following steps:

\begin{enumerate}
\item let $h=1,2,\ldots,(n-1)/2$;
\item let $q$ range over the nonzero odd integers satisfying $q^2h<n$, and
put $t=qh$ and $M=h(n-q^2h)$;
\item let $\nu$ range over the interval \eqref{eq:nu-interval};
\item compute $\Delta_\nu$ and the set of integral roots
\eqref{eq:mu-roots}; for each root form
$r=tz_0+\mu e_1+\nu e_2$;
\item retain $r$ only if $r_1\equiv r_3\pmod2$ and
$d=\tau^{-1}(r)$ is primitive, canonically oriented, and has odd third
coordinate;
\item form $F=-E+2qd$ and retain it only if $F\in\Tres$; write
$F=\eta F_0$, where $F_0=\pc([F])$ and $\eta\in\{\pm1\}$;
\item compare $d$ with the primitive oriented direction of $(E-F)/2$.
Retain the record
\[
 \bigl(\key(d),[F_0],\eta\bigr)
\]
precisely when $d$ has the smaller key.
\end{enumerate}

For each $h$, we place the surviving records in a finite set and sort them by
$(d_1,d_2,d_3)$.  Since $h$ increases in the outer loop, these blocks together
form the strict preference list of $[E]$.  The square roots and divisibility
tests use only integer arithmetic.  When the discriminant is zero, the two
signs in the quadratic formula contribute only one root.  The generator keeps one
sorted $h$-block at a time, returns its records, and then moves to the next
nonempty block.
\end{algorithm}

\begin{theorem}[Correctness of the preference-list algorithm]
\label{thm:incident-stream}
For every residual source orbit $[E]$, Algorithm~\ref{alg:incident-stream}
terminates and outputs every incident edge exactly once, in strictly increasing
order of its key.
\end{theorem}

\begin{proof}
We first check that every output is an incident edge.  A retained record
satisfies $Q(d)=h$ and $\Qip{E}{d}=qh$, and hence
\[
 Q(-E+2qd)=Q(E)-4q\Qip{E}{d}+4q^2Q(d)=n.
\]
The test $F\in\Tres$ gives the required congruence and exclusion conditions,
so the record defines a valid edge between the chosen representatives.  The
last comparison keeps the preferred sign of the target orbit, and $\eta$ is
the sign used in \eqref{eq:sign-lift}.

Conversely, take an incident edge and choose its preferred signed lift from $E$.
Lemma~\ref{lem:projective-half-bound} gives the stated range for
$h=Q(d)$.  Its vectors $p$ and $r=\tau(d)$ satisfy \eqref{eq:Pr}.  By
Lemma~\ref{lem:affine-coset}, $r$ has unique integral coordinates $\mu,\nu$, and
Lemma~\ref{lem:short-interval} shows that the procedure recovers them.  The
$\tau$-lattice, primitivity, orientation, parity, and residual tests all hold
for this lift, and the final comparison retains it.

It remains to check that no edge appears twice.  The values $h$ and $d$ are
intrinsic to a retained edge, while
$q=\Qip{E}{d}/Q(d)$ is then forced.  The affine coordinates $\mu,\nu$ are
unique, and \eqref{eq:mu-roots} is a set even when $\Delta_\nu=0$.  Hence no
record is emitted twice.
\end{proof}

\begin{example}[A complete preference list at $n=41$]
\label{ex:complete-row-41}
For $E=(-2,-5,2)$, direct enumeration of the residual target set gives four
target orbits.  Write their canonical representatives as
\[
\begin{aligned}
 F_1&=(-4,-3,0),&F_2&=(0,-3,4),\\
 F_3&=(-4,3,0), &F_4&=(0,-3,-4).
\end{aligned}
\]
Algorithm~\ref{alg:incident-stream} selects the following midpoint for each
orbit.  The last column records the $Q$-value of the
discarded midpoint direction obtained from the opposite sign.
\[
\begin{array}{c|c|c|c|c}
\text{orbit}&\eta&(E+\eta F_i)/2&\kappa([E],[F_i])&\text{other }Q\\ \hline
{[F_1]}&-1&(1,-1,1)&(5,1,-1,1)&36\\
{[F_2]}&-1&(-1,-1,-1)&(5,1,1,1)&36\\
{[F_3]}&-1&(1,-4,1)&(20,1,-4,1)&21\\
{[F_4]}& 1&(-1,-4,-1)&(20,1,4,1)&21
\end{array}
\]
The corresponding affine data $(h,q,\mu,\nu)$, in the same row order, are
\[
 (5,1,2,-1),\qquad
 (5,-1,2,1),\qquad
 (20,1,2,-4),\qquad
 (20,-1,2,4).
\]
For example, the third row is the full calculation in
Example~\ref{ex:lattice-41}; its opposite sign gives the direction
$(3,1,-1)$ of $Q$-value $21$.  In the second and fourth rows the displayed
midpoint is the negative of its canonically oriented primitive direction,
which is why the key has positive first coordinate.  Sorting the surviving
records gives the preference list
\[
 [F_1],\ [F_2],\ [F_3],\ [F_4]
\]
with keys in the order displayed.  Each of the four target orbits occurs
once.
\end{example}

\begin{remark}[A bound from lattice reduction]
\label{rem:size-reduction}
For a Lagrange--Gauss reduced basis,
\[
 |2\beta|\le\alpha\le\gamma.
\]
Equation \eqref{eq:gram-det} gives $\alpha\le2\sqrt{n/3}$.  The interval
\eqref{eq:nu-interval} therefore has a length bounded in terms of $n$ and $M$.
We do not need this bound for the correctness of the algorithm or for the
proposal counts below.
\end{remark}

\section{Deferred acceptance with generated preferences}
\label{sec:deferred-acceptance}

Proposition~\ref{prop:unique-stable} describes the matching by a global greedy
procedure.  To avoid sorting all edge keys at once, we use the
source-proposing Gale--Shapley algorithm \cite{GaleShapley} together with the
preference lists above.  We compute entries as they are requested and retain
the position reached in each list.  The order of the preferences remains
fixed throughout the algorithm.

Enumerate the finite residual orbit sets $\Sbar$ and $\Tbar$, and attach to
each source $s\in\Sbar$ the ordered generator of
Theorem~\ref{thm:incident-stream}.  A target compares two proposals by their
global edge rank.

\begin{algorithm}[Source-proposing deferred acceptance]\label{alg:lazy-da}
Initially every source is free and every target holds no proposal.
Place the sources in a first-in, first-out queue in lexicographic order of their
canonical representatives.  Repeatedly remove a free source $s$ from the
front of the queue and request the next record $(k,t,\eta)$ from its ordered
generator.
\begin{enumerate}
\item If $t$ holds no proposal, it holds $(s,k,\eta)$.
\item If $t$ holds $(s',k',\eta')$ and $k<k'$, then $t$ replaces
that record by $(s,k,\eta)$ and $s'$ is appended to the queue.
\item Otherwise $t$ rejects $s$, which is appended to the queue.
\end{enumerate}
The algorithm ends when the queue is empty.
\end{algorithm}

\begin{example}[Deferred acceptance on two sources]
\label{ex:DA-2x2}
Run Algorithm~\ref{alg:lazy-da} on the ranks in
Example~\ref{ex:ranked-2x2}, with initial queue $(s_1,s_2)$.  The proposals
are
\[
 s_1\longrightarrow t_1,
 \qquad
 s_2\longrightarrow t_1\quad\text{(rejected)},
 \qquad
 s_2\longrightarrow t_2.
\]
The first proposal is held by $t_1$.  The second has rank $2$, so $t_1$ keeps
its rank-$1$ proposal from $s_1$ and rejects $s_2$.  Source $s_2$ then
proposes to $t_2$.  The queue is empty and the output is
$\{(s_1,t_1),(s_2,t_2)\}$, the unique stable matching found by the greedy
argument.
\end{example}

\begin{theorem}[Correctness of deferred acceptance]\label{thm:DA-correct}
Algorithm~\ref{alg:lazy-da} terminates and returns the unique perfect stable
matching of Proposition~\ref{prop:unique-stable}.  The stored sign
$\eta$ determines the paired representatives through \eqref{eq:sign-lift}.
The inverse at a target is recovered by rerunning the same deterministic
matching algorithm; the sequence of earlier proposals need not be retained.
\end{theorem}

\begin{proof}
A source proposes in strict preference order and never proposes to the same
target twice.  We first check that a free source never exhausts its generator.
If a free source $s$ had proposed to all $m$ targets, then every target would
hold a proposal: once a target receives its first proposal, it remains
occupied thereafter.  The held proposals would come from $m$ distinct
sources, since a source can be held by at most one target and proposes again
only after becoming free.  None could come from the free source $s$, leaving
only $m-1$ possible sources, a contradiction.  Thus every request made by the
algorithm has a next record.  Since there are only $m^2$ possible proposals,
the algorithm terminates.

When the queue is empty, every source is held by a target.  There are $m$
sources and $m$ targets, and no target holds two sources, so the resulting
matching is perfect.

Suppose $(s,t)$ blocked the final matching.  Since $s$ prefers $t$ to its final
partner, $s$ proposed to $t$ earlier.  The target either rejected $s$
immediately or later displaced it for a proposal of smaller key.  A target's
held key only decreases, so its final edge is smaller than $(s,t)$, a
contradiction.  The output is stable and hence equals the unique stable
matching of Proposition~\ref{prop:unique-stable}.

For the inverse, let $F_0$ be the canonical representative of a matched target
orbit, and suppose its matched record is $(E_0,F_0,\eta)$ with $E_0$ canonical.
The lifted matching sends $E_0$ to $\eta F_0$.  Hence an actual target
$\rho F_0$, with $\rho\in\{\pm1\}$, has preimage $\rho\eta E_0$.
Lemma~\ref{lem:projective-key} makes $\eta$ intrinsic to the two orbits, and
rerunning the deterministic computation reconstructs the same matching and
the same sign.  This recovers the inverse at the given target without
retaining the first run's proposal transcript.
\end{proof}

\begin{example}[Recovering a residual preimage]
\label{ex:inverse-41}
Take the actual lifted target $F=(4,-3,0)$.  Its canonical orbit
representative is $F_0=(-4,3,0)$, so $F=\rho F_0$ with $\rho=-1$.
Rerunning the deterministic matching returns the record from
Example~\ref{ex:complete-row-41} with
\[
 E_0=(-2,-5,2),
 \qquad
 \eta=-1.
\]
The inverse sign formula gives
\[
 \rho\eta E_0=(-1)(-1)E_0=E_0.
\]
Finally, $\iota_B^{-1}(E_0)=(-2,-5,1)$.  Thus the inverse at the original
target is
\[
 (4,-3,0)\longmapsto(-2,-5,1).
\]
This inverse calculation requires rerunning the matching on all residual
orbits before choosing the sign.  It is not a formula depending only on the
coordinates of the given target.
\end{example}

\section{The map at \texorpdfstring{$n=41$}{n=41}}
\label{sec:example}

The quarter-turn calculation in Examples~\ref{ex:quarter-turn-41} and
\ref{ex:fiber-41} gives
\[
 (-4,-1,1)\longmapsto(2,-1,1)
 \quad\text{with even companion }(2,-1,2).
\]
We now work out the residual example
\[
 (-2,-5,1)\longmapsto(4,-3,0)
 \quad\text{with even companion }(4,-3,0).
\]
In each case, the even preimage is obtained by doubling the third coordinate
of the target.

\begin{example}[The residual matching for $n=41$]
\label{ex:41-residual}
The residual graph for $n=41$ has four source orbits $s_1,\ldots,s_4$ and
four target orbits $t_1,\ldots,t_4$.  We index them by the following
canonical lifted representatives:
\[
\begin{array}{ll}
 \pc(s_1)=(-2,-5,-2),&\pc(s_2)=(-2,-5,2),\\
 \pc(s_3)=(-2,5,-2), &\pc(s_4)=(-2,5,2),\\[2pt]
 \pc(t_1)=(-4,-3,0), &\pc(t_2)=(-4,3,0),\\
 \pc(t_3)=(0,-3,-4), &\pc(t_4)=(0,-3,4).
\end{array}
\]
In the next display, each row is in preference order and the bracketed
four-tuple is the corresponding edge key:
\[
\begin{aligned}
s_1&:\quad t_1[5,1,-1,-1],\ t_3[5,1,1,-1],
       \ t_2[20,1,-4,-1],\ t_4[20,1,4,-1],\\
s_2&:\quad t_1[5,1,-1,1],\ t_4[5,1,1,1],
       \ t_2[20,1,-4,1],\ t_3[20,1,4,1],\\
s_3&:\quad t_4[5,1,-1,-1],\ t_2[5,1,1,-1],
       \ t_3[20,1,-4,-1],\ t_1[20,1,4,-1],\\
s_4&:\quad t_3[5,1,-1,1],\ t_2[5,1,1,1],
       \ t_4[20,1,-4,1],\ t_1[20,1,4,1].
\end{aligned}
\]
The preference list in Example~\ref{ex:complete-row-41} is the
$s_2$ row here: its $F_1,F_2,F_3,F_4$ are respectively
$t_1,t_4,t_2,t_3$.
With the queue convention of Algorithm~\ref{alg:lazy-da}, the proposal
sequence is
\[
\begin{aligned}
s_1&\longrightarrow t_1,&
s_2&\longrightarrow t_1\quad\text{(rejected)},&
s_3&\longrightarrow t_4,\\
s_4&\longrightarrow t_3,&
s_2&\longrightarrow t_4\quad\text{(rejected)},&
s_2&\longrightarrow t_2.
\end{aligned}
\]
Thus the algorithm makes six proposals.  Its lifted residual pairs are
\[
\begin{aligned}
 (-2,-5,-2)&\longleftrightarrow(4,3,0),\\
 (-2,-5,2)&\longleftrightarrow(4,-3,0),\\
 (-2,5,-2)&\longleftrightarrow(0,-3,4),\\
 (-2,5,2)&\longleftrightarrow(0,-3,-4).
\end{aligned}
\]
For instance,
\[
 (-2,-5,1)\in\mathcal B_{41}^{\mathrm{odd}}
 \longmapsto
 (4,-3,0)\in\mathcal A_{41}.
\]
Thus
\[
\Phi_{41}^{-1}(4,-3,0)
 =\{(-2,-5,1),(4,-3,0)\}.
\]
\end{example}

\section{Parking functions and the cost of deferred acceptance}\label{sec:complexity}

A sequence $(a_1,\ldots,a_m)$ of positive integers is a \emph{parking
function} if its nondecreasing rearrangement $b_1\le\cdots\le b_m$ satisfies
$b_i\le i$ for $1\le i\le m$; see \cite{KonheimWeiss,StanleyParking}.
For $s\in\Sbar$, let $p(s)$ be the number of proposals made by $s$.
Equivalently, $p(s)$ is the rank of its final target in its preference list.
In the proof below, $M=\overline M_n$ denotes the final stable matching.

\begin{theorem}[The proposal counts form a parking function]
\label{thm:parking}
Order the sources $s_1,\ldots,s_m$ so that the keys of their final matched
edges are nondecreasing.  Then
\[
 p(s_i)\le i
 \qquad(1\le i\le m).
\]
Let $b_1\le\cdots\le b_m$ be the nondecreasing rearrangement of the proposal
counts.  Then $b_i\le i$, so the proposal-count sequence is a parking
function.  Moreover,
\begin{equation}\label{eq:proposal-bound}
 \sum_{s\in\Sbar}p(s)\le\frac{m(m+1)}2.
\end{equation}
\end{theorem}

\begin{proof}
Fix a source $s$, and let $t$ be a target that $s$ prefers to its final target.
Stability forces
\[
 \kappa(M^{-1}(t),t)<\kappa(s,t)<\kappa(s,M(s));
\]
otherwise $(s,t)$ would block.  Distinct preferred targets give distinct
matched edges of key strictly smaller than the matched edge of $s$.  Hence
\[
 p(s)-1\le
 \#\{e\in M:\kappa(e)<\kappa(s,M(s))\}.
\]
For $s=s_i$, the right-hand side is at most $i-1$, proving $p(s_i)\le i$.
For each $i$, the first $i$ sources in this ordering provide $i$ proposal
counts, all at most $i$.  Hence the $i$th smallest proposal count is at most
$i$, which is the parking-function condition.
Summing the inequalities $p(s_i)\le i$ gives \eqref{eq:proposal-bound}.
\end{proof}

\begin{example}[The parking function for $n=41$]
\label{ex:parking-41}
The residual run makes the six proposals
\[
 s_1\to t_1,
 \quad s_2\to t_1,
 \quad s_3\to t_4,
 \quad s_4\to t_3,
 \quad s_2\to t_4,
 \quad s_2\to t_2.
\]
Thus
\[
 p(s_1)=1,\qquad p(s_2)=3,\qquad p(s_3)=1,\qquad p(s_4)=1.
\]
Its nondecreasing rearrangement is $(1,1,1,3)$, and
\[
 1\le1,\qquad 1\le2,\qquad 1\le3,\qquad 3\le4.
\]
Thus it is a parking function.  The total number of proposals is $6$, while
Theorem~\ref{thm:parking} gives the bound
$4\cdot5/2=10$.  The algorithm requests six preference-list entries and
processes six proposals.  Exactly two of these proposals are compared with a
proposal already held by the target.  These are the operations counted in
Theorem~\ref{thm:complexity}.
\end{example}

\begin{proposition}[Sharpness]\label{prop:sharpness}
The bound \eqref{eq:proposal-bound} is sharp for the class of globally
ranked-pairs instances, even when every key class consists of a single edge.
\end{proposition}

\begin{proof}
Let the sides be $\{s_1,\ldots,s_m\}$ and $\{t_1,\ldots,t_m\}$, and give
$(s_i,t_j)$ the distinct label $(m+1)i+j$.  Every source orders the targets as
$t_1,t_2,\ldots,t_m$, and every target orders the sources as
$s_1,s_2,\ldots,s_m$.  In increasing-label greedy order, $(s_1,t_1)$ is
selected first; after its endpoints are removed, $(s_2,t_2)$ is selected,
and induction gives the diagonal stable matching.  The final target $t_i$ has
rank $i$ in the list of $s_i$, so $s_i$ makes exactly $i$ proposals.  The
total is $m(m+1)/2$.
\end{proof}

\begin{example}[Sharpness with three sources]
For $m=3$, all sources rank the targets as $t_1,t_2,t_3$, and all targets
rank the sources as $s_1,s_2,s_3$.  Starting with the queue
$(s_1,s_2,s_3)$, source $s_1$ is held by $t_1$.  Sources $s_2$ and $s_3$ are
then rejected by $t_1$; next $s_2$ is held by $t_2$, and $s_3$ is rejected by
$t_2$ before reaching $t_3$.  The proposal counts are therefore
\[
 p(s_1)=1,
 \qquad p(s_2)=2,
 \qquad p(s_3)=3,
\]
whose total $6=3\cdot4/2$ attains the bound.
\end{example}

\begin{theorem}[Cost of deferred acceptance]
\label{thm:complexity}
Let
\[
 N=\sum_{s\in\Sbar}p(s)
\]
be the number of proposals.  With preference lists supplied by the generators
of Theorem~\ref{thm:incident-stream}, the deferred-acceptance algorithm makes
$N$ requests for the next list entry, performs $O(N)$ key comparisons and
queue or assignment operations, and stores $O(m)$ auxiliary records apart
from those stored by the generators.
In particular,
\[
 N\le\frac{m(m+1)}2.
\]
The stated bound concerns the deferred-acceptance phase only.  It does not
include the arithmetic work or storage internal to the generators, the
enumeration of the residual orbits, or integer bit complexity.
\end{theorem}

\begin{proof}
Processing a proposal requires one key comparison if the target already
holds a proposal, and none otherwise.  It also requires a constant number of
queue and matching-array updates.  Apart from the generators, the algorithm
stores the status of each source, one held record per target, and the queue
of free sources.  These require $O(m)$ records.  The bound on $N$ is
Theorem~\ref{thm:parking}.
\end{proof}

The theorem does not bound the total cost of evaluating $\Phi_n$: producing
the next preference-list entry may require further arithmetic.  Computing an
inverse also requires another run of the matching algorithm, whose cost is
not included above.

\begin{remark}[On possible subquadratic bounds]\label{rem:subquadratic}
Proposition~\ref{prop:sharpness} shows that an $O(m\log m)$ worst-case bound
on the number of proposals cannot follow from the globally-ranked-pairs
property alone.  Such a bound for the Tunnell-form instances considered here
would require additional
arithmetic information about the distribution of their midpoint directions.
\end{remark}

\section{Concluding remarks}\label{sec:audit}

We have constructed the odd-branch bijection using three quarter-turns and a
stable matching on the remaining antipodal orbits.  Primitive midpoint
directions determine the preferences, and the common ranking of edges gives
both uniqueness and the parking-function bound.  The examples at $n=41$
show the two cases: an odd preimage found by a quarter-turn and one found by
the residual matching.  In both cases, doubling the third coordinate of the
target gives the even preimage.

The preference lists can be generated without storing the complete
$m\times m$ edge table.  It remains to determine whether the arithmetic
distribution of midpoint directions gives a better complexity bound for
these representation sets.  A separate question is whether the residual
partner can be found directly from one representation, without computing
the global matching.

\section*{Acknowledgments}

This work grew out of the 2026 Research in Industrial Projects for Students
(RIPS) program at the Institute for Pure \& Applied Mathematics (IPAM) at
UCLA.  The authors thank IPAM for providing a collaborative research
environment.

\section*{Funding}

This work was supported by OpenAI through its sponsorship of the RIPS project
and its provision of access to the computational tools used in the project.

\section*{Declaration of generative AI and AI-assisted technologies in the
manuscript preparation process}

During the mathematical development and preparation of this work, the authors
used OpenAI's ChatGPT and Codex through an AI-assisted solver--referee pipeline
for symbolic experimentation, literature searches, proof organization, and
language revision.

\bibliographystyle{amsplain}
\bibliography{references}

\end{document}
